\subsection{Pointed homotopies}

Let X and Y be topological spaces and let x be a point of X. A homotopy \(\sigma\colon X \times I \to Y\) is said to be pointed at x if it is fixed on \(\{x\}\), that is, if the map \(t \mapsto \sigma(x, t)\) from I to Y is constant. The juxtaposition of two homotopies pointed at x is pointed at x. The origin and the term of a homotopy pointed at x take the same value y at x; they are therefore pointed continuous maps from \((X, x)\) to \((Y, y)\).

Let \((X,x)\) and \((Y,y)\) be pointed topological spaces (I, p. 120, definition 1). The relation “f is connected to g by a homotopy pointed at x” is an equivalence relation on the set \(\mathcal{C}((X,x);(Y,y))\), called the pointed homotopy relation. The quotient set of \(\mathcal{C}((X,x);(Y,y))\) by this equivalence relation is denoted \([(X,x);(Y,y)]\). Its elements are called pointed homotopy classes of pointed continuous maps from \((X,x)\) to \((Y,y)\).

\pandocbounded{\includegraphics[keepaspectratio]{images/part002_b3045f29.jpg}}

The set \(\mathcal{C}((\mathrm{X},x);(\mathrm{Y},y))\) is a subset of \(\mathcal{C}(\mathrm{X};\mathrm{Y})\). Note that pointed homotopy is an equivalence relation finer than the relation induced by homotopy on \(\mathcal{C}(\mathrm{X};\mathrm{Y})\). It is in general a strictly finer relation (see III, p. 321, exerc. 1 and III, p. 234, example 3).

Let also \((Z,z)\) be a pointed topological space. Let f and \(f'\) be pointed continuous maps from \((X,x)\) to \((Y,y)\), and g and \(g'\) be pointed continuous maps from \((Y,y)\) to \((Z,z)\). If f and \(f'\) are related by a pointed homotopy \(\sigma\) at x, and if g and \(g'\) are related by a pointed homotopy \(\tau\) pointed at y, then \(g \circ f\) and \(g' \circ f'\) are connected by the homotopy \(\theta: X \times I \to Z\), \((u,t) \mapsto \tau(\sigma(u,t),t)\), which is pointed at x. As above, this allows one to define the composition map of \([(X,x);(Y,y)] \times [(Y,y);(Z,z)]\) to \([(X,x);(Z,z)]\). We leave it to the reader to formulate and prove for pointed maps and pointed homotopies the statements analogous to Prop. 3 and its corollary.

